\documentclass[12pt, a4paper]{article}
\setlength{\oddsidemargin}{0.5cm}
\setlength{\evensidemargin}{0.5cm}
\setlength{\topmargin}{-1cm}
\setlength{\leftmargin}{0.5cm}
\setlength{\rightmargin}{0.5cm}
\setlength{\textheight}{24.00cm} 
\setlength{\textwidth}{15.00cm}
\parindent 0pt
\parskip 5pt
\pagestyle{plain}
\usepackage{graphicx}
\usepackage{url}
\usepackage{titling}
\usepackage{float}
\setlength{\droptitle}{-3cm}
% \postdate{\par\vspace{-1cm}}
\posttitle{\par\vspace{-1cm}}

\title{Research Proposal}
\date{March 2026}

\newcommand{\namelistlabel}[1]{\mbox{#1}\hfil}
\newenvironment{namelist}[1]{%1
\begin{list}{}
    {
        \let\makelabel\namelistlabel
        \settowidth{\labelwidth}{#1}
        \setlength{\leftmargin}{1.1\labelwidth}
    }
  }{%1
\end{list}}

\begin{document}
\maketitle

\begin{namelist}{xxxxxxxxxxxx}
\item[{\bf Title:}]
	Training-Free Agentic RAG Framework for Multi-Page Visually Rich Document Understanding
\item[{\bf Author:}]
	Lewei Xu
\item[{\bf Supervisor:}]
	Dr. Yihao Ding, Dr. Daochang Liu
\item[{\bf Degree:}]
	BACS (Honours)
\end{namelist}

\section*{Background}
% Emphasise the signifance of multi page document understanding. (An example can be given). Explain why multipage document is difficult: long context, long logical connection, multi-hop

% General Development from single page, unimodality understanding to multipage, multipage document understanding

Modern organizations rely heavily on long-form documents such as contracts, invoices, financial reports, and scientific articles, yet extracting specific information from these documents remains a time-consuming and largely manual process. This challenge arises because such documents contain not only text, but also visual elements such as tables, figures, and graphs, and are frequently scanned images rather than well-formatted PDFs with embedded text \cite{DocExtractNet}. Consequently, a model designed to answer questions about these documents must account for both the spatial layout of components within a single page and the logical relationships between components spread across multiple pages. This is a task that Document Visual Question Answering (DocVQA) models have historically struggled with \cite{DocVQA}.

Early DocVQA research focused exclusively on single-page documents, where the task was already non-trivial due to the need for joint understanding of text, layout, and visual content. This line of research has largely reached maturity with the introduction of layout-aware models such as LayoutLLM \cite{LayoutLLM} and the subsequent emergence of multimodal large language models (MLLMs) such as Qwen3-VL \cite{Qwen3-VL}, which now achieves near human-level performance on single-page benchmarks \cite{rrcsingle}.

However, real-world documents are rarely confined to a single page. As documents grow longer, the complexity of the task increases substantially: single-page models struggle to fit entire documents within their context window \cite{MP-DocVQA}, and tend to treat each page as an isolated entity, failing to capture logical connections and multi-hop logical connections that span multiple pages \cite{GRAM}. For example, answering a question such as ``What is the year-on-year change in total revenue?'' may require locating financial figures on separate pages and reasoning over their relationship, a capability that single-page models fundamentally lack.

Multi-Page Document Visual Question Answering (MP-DocVQA) has emerged to address the limitations of Single-Page DocVQA. Early approaches such as Hi-VT5 \cite{MP-DocVQA} introduced hierarchical aggregation to consolidate information across pages, while more recent methods have incorporated Retrieval-Augmented Generation (RAG) \cite{RAG-DocVQA} to selectively retrieve relevant pages before generating an answer. Despite these advances, MP-DocVQA models still fall well short of human-level performance \cite{rrcmulti}.

\section*{Related Work}

\subsection*{MLLM-Based Approaches}
Recent multimodal large language model (MLLM) based approaches to Multi-Page DocVQA process entire documents within extended context windows. DocLayLLM \cite{DocLayLLM} encodes both textual and layout information within a unified MLLM architecture, achieving strong zero-shot performance across document understanding tasks. DocoPilot \cite{Docopilot} takes a complementary approach, processing documents natively through engineering optimisations such as data packing and Ring Attention. Despite these advances, both models degrade as document length approaches or exceeds their context window, and their single-pass designs provide no mechanism for iterative reasoning or targeted page inspection.

\subsection*{RAG-Based and Agentic Approaches}
RAG addresses the context window limitations of single-pass models by retrieving relevant pages prior to generation. RAG-DocVQA demonstrates gains of up to +22.5 ANLS on MP-DocVQA over the Hi-VT5 baseline \cite{RAG-DocVQA}, while M3DocRAG \cite{M3DocRAG} extends this to multi-modal retrieval. MDocAgent \cite{MDocAgent} and ViDoRAG \cite{ViDoRAG} further incorporate multi-agent and iterative retrieval workflows. However, all such approaches rely on fixed retrieval pipelines and require task-specific training, limiting their adaptability to unseen document types and complex multi-hop questions.

\subsection*{Research Gaps}
Three key limitations persist in current approaches to multi-page VRDU. First, all existing approaches require task-specific training on labelled data, restricting generalisation to new document domains. Second, fixed RAG retrieval workflows lack a mechanism to explore alternative retrieval trajectories when initial retrieval fails. Third, cross-page fusion of multi-modal evidence across tables, figures, and text in long documents remains an open problem.

\section*{Research Aims}

\subsection*{Aim 1: Develop a Training-Free Knowledge Graph-Augmented RAG Pipeline for Multi-Page DocVQA}
The first aim is to develop a training-free RAG pipeline augmented with a structured knowledge graph constructed directly from the document being queried. This improves multi-hop retrieval across multi-page visually rich documents without requiring task-specific fine-tuning, directly addressing the first and third research gaps identified above.

\subsection*{Aim 2: Improve Retrieval Performance via Agentic Planning and a Specialised Tool Set}
The second aim is to augment the RAG pipeline developed in Aim 1 with an agentic planning framework that adaptively selects and sequences a document-specialised tool set to iteratively retrieve and process evidence from the knowledge graph. This directly addresses the second research gap by replacing fixed, single-pass retrieval with a principled search mechanism capable of exploring alternative retrieval trajectories when initial retrieval fails.
 
\section*{Method}
This project will be carried out in 4 phases across 2 semesters between March and November 2026. The timeframes of the 4 phases are expected to overlap significantly.

\subsection*{Phase 1: Literature Review}
\subsubsection*{Expected Duration: March - May (10 weeks)}
All relevant papers will be aggregated into a structured spreadsheet cataloguing each work's key technique, year, benchmark performance, and relevance to the proposed research. This will then be consolidated into a formal literature review covering four areas:
\begin{enumerate}
    \item Single-page DocVQA models and their limitations on multi-page benchmarks.
    \item Multi-page DocVQA benchmarks and dataset selection (MP-DocVQA, DUDE, MMDocVQA, MMLongBench-Doc).
    \item Existing RAG-based and agentic approaches (M3DocRAG, MDocAgent, ViDoRAG).
    \item In-depth study of the ToolTree framework \cite{ToolTree} to identify modifications required for the multi-page DocVQA setting.
\end{enumerate}

\subsection*{Phase 2: Design and Development}
\subsubsection*{Expected Duration: Late April - August (16 Weeks)}

\textbf{Training-Free RAG + KG (Aim 1).} A structured knowledge graph will be constructed from target documents using off-the-shelf layout parsers (DocLayNet, PPStructure) to segment pages into semantic units (paragraphs, tables, figures, captions etc.). Three types of edges will be established between nodes: explicit cross-reference links detected via rule-based text parsing, structural hierarchy links derived from heading structure, and semantic similarity links computed using pre-trained text and visual embeddings (BGE, ColPali). 

\textbf{Agentic Planning (Aim 2).} A document-specialised tool set will be designed and implemented, including visual page retrieval, OCR, table and figure parsing, and cross-page evidence aggregation. Each tool will be accompanied by a structured tool card describing its input/output schema to enable training-free tool selection. The planning strategy may extend ToolTree's \cite{ToolTree} MCTS framework, adapting the pre-evaluation scoring ($r_{pre}$) and post-evaluation scoring ($r_{post}$) to a multi-page setting.

\subsection*{Phase 3: Experimental Evaluation and Analysis}
\subsubsection*{Expected Duration: July - September (12 Weeks)}
The proposed framework will be evaluated on the benchmarks selected in Phase 1, comparing against existing RAG-based, MLLM-based, and agentic baselines using ANLS as the primary metric. Ablation studies will be conducted to isolate the contribution of each component: the knowledge base design, tool set, and agentic planning strategy.  Qualitative error analysis will be used to identify remaining weaknesses and failure modes.

\subsection*{Phase 4: Dissertation and Presentation}
\subsubsection*{Expected Duration: Late August - November (10 Weeks)}
All experimental results, ablation findings, and error analyses will be consolidated into the final dissertation. A seminar presentation and accompanying script will be prepared to communicate the key contributions and findings, and a live demonstration website may be developed to showcase the proposed framework in an interactive setting.

\section*{Software and Hardware Requirements}
This project will be completed on the following system: Intel 265k (20 cores), 64gb DDR5 RAM, NVIDIA RTX 5070 12gb VRAM with CUDA. 

The development environment will be Windows Subsystem for Linux (WSL), with all implementation carried out in Python using open-source libraries and frameworks. 

As proposed project is training free, the hardware specified above should be sufficient for the project, but should any training of the framework components be required, access to dedicated external GPU's may be necessary.

\begin{figure}
    \centering
    \includegraphics[width=1\linewidth]{chart.png}
    \label{fig:placeholder}
\end{figure}

\newpage

\section*{Statement of AI Use}
LLMs were used in this document to help clean up grammar, wording, and formatting. This document was completed in overleaf, original text can be retrieved on request via overleaf versioning.

\begin{thebibliography}{9}

\bibitem{DocExtractNet}
Z. Yan, Z. Ye, J. Ge, J. Qin, J. Liu, Y. Cheng, and C. Gurrin, ``DocExtractNet: A novel framework for enhanced information extraction from business documents,'' \textit{Information Processing \& Management}, vol. 60, no. 3, Article 104046, 2024. doi: 10.1016/j.ipm.2024.104046

\bibitem{DocVQA}
M. Mathew, D. Karatzas, and C. V. Jawahar, ``DocVQA: A dataset for VQA on document images,'' in \textit{Proc. IEEE/CVF Winter Conf. on Applications of Computer Vision (WACV)}, 2021. arXiv:2007.00398

\bibitem{LayoutLLM}
C. Luo, Y. Shen, Z. Zhu, Q. Zheng, Z. Yu, and C. Yao, ``LayoutLLM: Layout instruction tuning with large language models for document understanding,'' in \textit{Proc. IEEE/CVF Conf. on Computer Vision and Pattern Recognition (CVPR)}, 2024. arXiv:2404.05225

\bibitem{Qwen3-VL}
S. Bai \textit{et al.}, ``Qwen3-VL technical report,'' \textit{arXiv preprint arXiv:2511.21631}, 2025.

\bibitem{rrcsingle}
Robust Reading Competition, ``Task 1: Single-page document visual question answering,'' [Online]. Available: \url{https://rrc.cvc.uab.es/?ch=17&com=evaluation&task=1} [Accessed: Mar. 2026]

\bibitem{MP-DocVQA}
R. Tito, D. Karatzas, and E. Valveny, ``Hierarchical multimodal transformers for multipage DocVQA,'' \textit{Pattern Recognition}, vol. 144, Article 109834, 2023. doi: 10.1016/j.patcog.2023.109834

\bibitem{GRAM}
T. Blau, S. Fogel, R. Ronen, A. Golts, R. Ganz, E. Ben Avraham, A. Aberdam, S. Tsiper, and R. Litman, ``GRAM: Global reasoning for multi-page VQA,'' in \textit{Proc. IEEE/CVF Conf. on Computer Vision and Pattern Recognition (CVPR)}, 2024. arXiv:2401.03411

\bibitem{RAG-DocVQA}
E. López \textit{et al.}, ``Enhancing document VQA models via retrieval-augmented generation,'' in \textit{Proc. Int. Conf. on Document Analysis and Recognition (ICDAR)}, 2025. arXiv:2508.18984

\bibitem{rrcmulti}
Robust Reading Competition, ``Task 4: Multi-page document visual question answering,'' [Online]. Available: \url{https://rrc.cvc.uab.es/?ch=17&com=evaluation&task=4} [Accessed: Mar. 2026]

\bibitem{DocLayLLM}
W. Liao, J. Wang, H. Li, C. Wang, J. Huang, and L. Jin, ``DocLayLLM: An efficient and effective multi-modal extension of large language models for text-rich document understanding,'' \textit{arXiv preprint arXiv:2408.15045}, 2024.

\bibitem{Docopilot}
Y. Duan, Z. Chen, Y. Hu, W. Wang, S. Ye, B. Shi, L. Lu, Q. Hou, T. Lu, H. Li, J. Dai, and W. Wang, ``Docopilot: Improving multimodal models for document-level understanding,'' in \textit{Proc. IEEE/CVF Conf. on Computer Vision and Pattern Recognition (CVPR)}, pp. 4026--4037, 2025. arXiv:2507.14675

\bibitem{M3DocRAG}
J. Cho, D. Mahata, O. Irsoy, Y. He, and M. Bansal, ``M3DocRAG: Multi-modal retrieval is what you need for multi-page multi-document understanding,'' \textit{arXiv preprint arXiv:2411.04952}, 2024.

\bibitem{MDocAgent}
S. Han, P. Xia, R. Zhang, T. Sun, Y. Li, H. Zhu, and H. Yao, ``MDocAgent: A multi-modal multi-agent framework for document understanding,'' \textit{arXiv preprint arXiv:2503.13964}, 2025.

\bibitem{ViDoRAG}
Q. Wang, R. Ding, Z. Chen, W. Wu, S. Wang, P. Xie, and F. Zhao, ``ViDoRAG: Visual document retrieval-augmented generation via dynamic iterative reasoning agents,'' \textit{arXiv preprint arXiv:2502.18017}, 2025.

\bibitem{ToolTree}
S. Yang, S. C. Han, Y. Ding, S. Wang, and E. Hovy, ``ToolTree: Efficient LLM agent tool planning via dual-feedback Monte Carlo tree search and bidirectional pruning,'' in \textit{Proc. 14th Int. Conf. on Learning Representations (ICLR)}, 2026.

\end{thebibliography}
\end{document}

